% !TEX program = pdflatex
% ==========================================================
% CLASE BEAMER TOP
% Fundamentos de Matemática Básica - Agronomía
% Tema: Potencias
% Universidad Santo Tomás - Talca
% ==========================================================

\newif\ifhandout
%\handoutfalse
\handouttrue

\ifhandout
\documentclass[aspectratio=169,10pt,handout]{beamer}
\else
\documentclass[aspectratio=169,10pt]{beamer}
\fi

% ----------------------------------------------------------
% PAQUETES
% ----------------------------------------------------------
\usepackage[spanish,es-noshorthands]{babel}
\usepackage[T1]{fontenc}
\usepackage[utf8]{inputenc}
\usepackage{lmodern}
\usepackage{amsmath,amsfonts,amssymb}
\usepackage{ragged2e}
\usepackage{enumitem}
\usepackage{multicol}
\usepackage{array}
\usepackage{booktabs}
\usepackage{tikz}
\usetikzlibrary{positioning,calc}
\usepackage{fontawesome5}

% ----------------------------------------------------------
% COLORES
% ----------------------------------------------------------
\definecolor{USTBlue}{RGB}{0,70,127}
\definecolor{USTBlueDark}{RGB}{0,45,85}
\definecolor{USTTeal}{RGB}{0,153,117}
\definecolor{USTGray}{RGB}{120,120,120}
\definecolor{USTSoft}{RGB}{248,249,250}
\definecolor{USTText}{RGB}{35,40,45}
\definecolor{USTGold}{RGB}{196,154,70}
\definecolor{USTLightBlue}{RGB}{227,239,248}
\definecolor{USTLightTeal}{RGB}{226,245,240}
\definecolor{USTRedSoft}{RGB}{252,235,235}
\definecolor{USTLine}{RGB}{220,226,232}

% ----------------------------------------------------------
% TEMA
% ----------------------------------------------------------
\usetheme{Madrid}
\usecolortheme[named=USTBlue]{structure}
\setbeamertemplate{navigation symbols}{}

% ----------------------------------------------------------
% AJUSTES VISUALES
% ----------------------------------------------------------
\setlist[itemize]{itemsep=0.45em}
\setlist[enumerate]{itemsep=0.45em}

\setbeamercolor{normal text}{fg=USTText,bg=white}
\setbeamercolor{title}{fg=USTBlueDark}
\setbeamercolor{subtitle}{fg=USTTeal}
\setbeamercolor{frametitle}{fg=white,bg=USTBlueDark}

\setbeamercolor{block title}{bg=USTLightBlue,fg=USTBlueDark}
\setbeamercolor{block body}{bg=USTSoft,fg=black}

\setbeamercolor{block title example}{bg=USTLightTeal,fg=USTBlueDark}
\setbeamercolor{block body example}{bg=USTSoft,fg=black}

\setbeamercolor{block title alerted}{bg=USTRedSoft,fg=black}
\setbeamercolor{block body alerted}{bg=red!2,fg=black}

% ----------------------------------------------------------
% FRAMETITLE
% ----------------------------------------------------------
\setbeamertemplate{frametitle}{
	\vspace*{-0.03cm}
	\begin{beamercolorbox}[wd=\paperwidth,ht=0.92cm,dp=0.18cm,leftskip=0.45cm,rightskip=0.25cm]{frametitle}
		{\large\bfseries \insertframetitle}
	\end{beamercolorbox}
}

% ----------------------------------------------------------
% PIE DE PÁGINA
% ----------------------------------------------------------
\setbeamercolor{author in head/foot}{bg=USTSoft,fg=USTGray}
\setbeamercolor{date in head/foot}{bg=USTSoft,fg=USTGray}
\setbeamercolor{title in head/foot}{bg=USTSoft,fg=USTGray}

\setbeamertemplate{footline}{%
	\leavevmode%
	\hbox{%
		\begin{beamercolorbox}[wd=.82\paperwidth,ht=2.8ex,dp=1ex,leftskip=0.4cm]{author in head/foot}%
			\scriptsize
			\textbf{Fundamentos de Matemática Básica}
			\quad|\quad
			Agronomía
			\quad|\quad
			Universidad Santo Tomás Talca
		\end{beamercolorbox}%
		\begin{beamercolorbox}[wd=.18\paperwidth,ht=2.8ex,dp=1ex,rightskip=0.35cm plus1fil]{date in head/foot}%
			\scriptsize \insertframenumber/\inserttotalframenumber
		\end{beamercolorbox}%
	}%
}

% ----------------------------------------------------------
% COMANDOS
% ----------------------------------------------------------
\newcommand{\destacar}[1]{\textcolor{USTBlueDark}{\textbf{#1}}}
\newcommand{\alerta}[1]{\textcolor{red!70!black}{\textbf{#1}}}
\newcommand{\pp}{\pause}

\newcommand{\iconoInicio}{\faPlayCircle}
\newcommand{\iconoDesarrollo}{\faBrain}
\newcommand{\iconoPractica}{\faPenSquare}
\newcommand{\iconoAgro}{\faSeedling}
\newcommand{\iconoCierre}{\faCheckCircle}
\newcommand{\iconoIdea}{\faLightbulb}
\newcommand{\iconoError}{\faExclamationTriangle}
\newcommand{\iconoTiempo}{\faClock}
\newcommand{\iconoMeta}{\faBullseye}

\newenvironment{metaBox}
{\begin{block}{\iconoMeta\ \ Meta de aprendizaje}}
	{\end{block}}

\newenvironment{ideaBox}
{\begin{exampleblock}{\iconoIdea\ \ Idea clave}}
	{\end{exampleblock}}

% ----------------------------------------------------------
% PORTADA TOP
% ----------------------------------------------------------
\newcommand{\PortadaProfesionalUST}{%
	\begin{frame}[plain]
		\begin{tikzpicture}[remember picture,overlay]
			
			\fill[white] (current page.south west) rectangle (current page.north east);
			\fill[USTBlueDark] (current page.north west) rectangle ([yshift=-1.45cm]current page.north east);
			\fill[USTTeal] ([yshift=-1.45cm]current page.north west) rectangle ([yshift=-1.62cm]current page.north east);
			\fill[USTSoft] (current page.south west) rectangle ([xshift=2.25cm]current page.north west);
			
			\fill[USTBlue!7] ([xshift=1.10cm,yshift=-2.0cm]current page.north west) circle (0.78cm);
			\fill[USTTeal!10] ([xshift=1.70cm,yshift=-4.2cm]current page.north west) circle (0.45cm);
			\fill[USTGold!14] ([xshift=-1.2cm,yshift=1cm]current page.south east) circle (1.08cm);
			
			\fill[black!7,rounded corners=8pt]
			([xshift=-5.1cm,yshift=-3.15cm]current page.center) rectangle
			([xshift=5.0cm,yshift=2.75cm]current page.center);
			
			\fill[white,rounded corners=8pt]
			([xshift=-5.25cm,yshift=-3.00cm]current page.center) rectangle
			([xshift=4.85cm,yshift=2.90cm]current page.center);
			
			\fill[USTTeal] ([xshift=-5.25cm,yshift=-3.00cm]current page.center) rectangle
			([xshift=-4.95cm,yshift=2.90cm]current page.center);
			
			\node[anchor=north east,text=white]
			at ([xshift=-0.7cm,yshift=-0.52cm]current page.north east)
			{\small\bfseries Universidad Santo Tomás \quad Ciencias Básicas Talca};
			
			\node[align=center] at ([xshift=0.10cm,yshift=-0.05cm]current page.center){%
				\begin{minipage}{0.68\paperwidth}
					\centering
					
					{\Large\bfseries\color{USTText}\insertauthor\par}
					
					\vspace{0.75cm}
					
					{\fontsize{24}{28}\selectfont\bfseries\color{USTBlueDark}\inserttitle\par}
					
					\vspace{0.28cm}
					
					{\large\bfseries\color{USTTeal}\insertsubtitle\par}
					
					\vspace{0.58cm}
					
					{\normalsize\color{USTGray}\faSeedling\ Agronomía \quad \faUniversity\ Universidad Santo Tomás\par}
					
					\vspace{0.65cm}
					
					{\small\color{USTGray}\insertdate\par}
				\end{minipage}
			};
			
			\draw[USTGray!35,line width=0.35pt]
			([xshift=0.9cm,yshift=0.72cm]current page.south west) --
			([xshift=-0.9cm,yshift=0.72cm]current page.south east);
			
			\node[anchor=south,text=USTGray]
			at ([yshift=0.2cm]current page.south)
			{\scriptsize
				\textbf{Fundamentos de Matemática Básica}
				\;|\;
				Clase TOP
				\;|\;
				Potencias};
		\end{tikzpicture}
	\end{frame}
}

% ----------------------------------------------------------
% FRAME DE SECCIÓN
% ----------------------------------------------------------
\newcommand{\SectionFrame}[3]{%
	\begin{frame}[plain]
		\begin{tikzpicture}[remember picture,overlay]
			\fill[USTBlueDark] (current page.south west) rectangle (current page.north east);
			\fill[USTTeal] (current page.south west) rectangle ([yshift=0.23cm]current page.north east);
			\fill[white,opacity=0.08] ([xshift=-1cm,yshift=1cm]current page.south east) circle (1.8cm);
			\fill[white,opacity=0.06] ([xshift=1.3cm,yshift=-0.8cm]current page.north west) circle (1.25cm);
			
			\node[anchor=center,text=white] at ([yshift=0.9cm]current page.center)
			{{\fontsize{34}{36}\selectfont #1}};
			
			\node[anchor=center,text=white] at (current page.center)
			{{\LARGE\bfseries #2}};
			
			\node[anchor=center,text=white!90] at ([yshift=-1.0cm]current page.center)
			{{\large #3}};
		\end{tikzpicture}
	\end{frame}
}

% ----------------------------------------------------------
% DATOS
% ----------------------------------------------------------
\title{Fundamentos de Matemática Básica}
\subtitle{Clase 2: Potencias}
\author{Prof. Luis González Alcaino}
\date{Universidad Santo Tomás - Agronomía 2026}

% ==========================================================
% DOCUMENTO
% ==========================================================
\begin{document}
	
	% ----------------------------------------------------------
	\begin{frame}
		\titlepage
	\end{frame}
	
		\begin{metaBox}
			Aplicar el concepto de potencia y sus propiedades en ejercicios matemáticos y situaciones contextualizadas en Agronomía.
		\end{metaBox}
		
		\pp
		
		\begin{block}{\faListOl\ \ Ruta de la sesión}
			\begin{enumerate}
				\item Inicio
				\item Desarrollo conceptual
				\item Práctica guiada
				\item Aplicación agronómica
				\item Práctica autónoma
				\item Cierre
			\end{enumerate}
		\end{block}
	\end{frame}
	
	% ----------------------------------------------------------
	\begin{frame}{\iconoTiempo\ \ Distribución del tiempo}
		\centering
		\begin{tabular}{>{\raggedright\arraybackslash}p{4.4cm} >{\centering\arraybackslash}p{2.2cm} >{\raggedright\arraybackslash}p{4.6cm}}
			\toprule
			\textbf{Momento} & \textbf{Tiempo} & \textbf{Propósito} \\
			\midrule
			Inicio & 10 min & Activar conocimientos previos \\
			Desarrollo & 25 min & Comprender concepto y propiedades \\
			Práctica guiada & 15 min & Resolver paso a paso \\
			Aplicación agronómica & 10 min & Conectar con el contexto profesional \\
			Práctica autónoma & 15 min & Consolidar individualmente \\
			Cierre & 5 min & Sintetizar y reflexionar \\
			\bottomrule
		\end{tabular}
	\end{frame}
	
	% ----------------------------------------------------------
	\SectionFrame{\iconoInicio}{Inicio}{Activación y conexión inicial}
	
	% ----------------------------------------------------------
	\begin{frame}{\iconoInicio\ \ Activación inicial}
		\begin{block}{Pensemos juntos}
			\begin{itemize}
				\item<1-> ¿Cómo escribirías \(2\cdot2\cdot2\cdot2\) de manera abreviada?
				\item<2-> ¿Qué representa el exponente?
				\item<3-> ¿Dónde aparecen multiplicaciones repetidas en Agronomía?
			\end{itemize}
		\end{block}
		
		\only<4>{
			\begin{ideaBox}
				Las potencias permiten expresar crecimiento, repetición y escala de forma compacta.
			\end{ideaBox}
		}
	\end{frame}
	
	% ----------------------------------------------------------
	\begin{frame}{\iconoInicio\ \ Motivación}
		\begin{columns}[T,totalwidth=\textwidth]
			\begin{column}{0.48\textwidth}
				\begin{block}{Usos en Agronomía}
					\begin{itemize}
						\item Densidad de siembra
						\item Crecimiento bacteriano
						\item Micronutrientes
						\item Notación científica
					\end{itemize}
				\end{block}
			\end{column}
			\begin{column}{0.48\textwidth}
				\begin{exampleblock}{Ejemplo}
					\[
					0{,}00045\ \text{kg}=4{,}5\times10^{-4}\ \text{kg}
					\]
					Esta escritura permite calcular y comunicar con mayor claridad.
				\end{exampleblock}
			\end{column}
		\end{columns}
	\end{frame}
	
	% ----------------------------------------------------------
	\SectionFrame{\iconoDesarrollo}{Desarrollo}{Concepto y propiedades de las potencias}
	
	% ----------------------------------------------------------
	\begin{frame}{\iconoDesarrollo\ \ Concepto de potencia}
		\begin{block}{Definición}
			\[
			a^n=\underbrace{a\cdot a\cdot a\cdots a}_{n\text{ veces}}
			\]
		\end{block}
		
		\pp
		
		\begin{columns}[T,totalwidth=\textwidth]
			\begin{column}{0.47\textwidth}
				\begin{block}{Base}
					Número que se repite como factor.
				\end{block}
			\end{column}
			\begin{column}{0.47\textwidth}
				\begin{exampleblock}{Exponente}
					Cantidad de veces que se repite la base.
				\end{exampleblock}
			\end{column}
		\end{columns}
		
		\pp
		
		\begin{exampleblock}{Ejemplo}
			\[
			5^3=5\cdot5\cdot5=125
			\]
		\end{exampleblock}
	\end{frame}
	
	% ----------------------------------------------------------
	\begin{frame}{\iconoDesarrollo\ \ Propiedad 1: producto de igual base}
		\[
		a^m\cdot a^n=a^{m+n}
		\]
		
		\pp
		
		\begin{exampleblock}{Ejemplo}
			\[
			2^3\cdot2^4=2^{3+4}=2^7=128
			\]
		\end{exampleblock}
		
		\pp
		
		\begin{ideaBox}
			Misma base \(\rightarrow\) se suman exponentes.
		\end{ideaBox}
	\end{frame}
	
	% ----------------------------------------------------------
	\begin{frame}{\iconoDesarrollo\ \ Propiedad 2: cociente de igual base}
		\[
		\frac{a^m}{a^n}=a^{m-n}, \qquad a\neq0
		\]
		
		\pp
		
		\begin{exampleblock}{Ejemplo}
			\[
			\frac{5^6}{5^2}=5^{6-2}=5^4=625
			\]
		\end{exampleblock}
		
		\pp
		
		\begin{ideaBox}
			Misma base \(\rightarrow\) se restan exponentes.
		\end{ideaBox}
	\end{frame}
	
	% ----------------------------------------------------------
	\begin{frame}{\iconoDesarrollo\ \ Propiedad 3: potencia de una potencia}
		\[
		(a^m)^n=a^{m\cdot n}
		\]
		
		\pp
		
		\begin{exampleblock}{Ejemplo}
			\[
			(3^2)^4=3^{2\cdot4}=3^8=6561
			\]
		\end{exampleblock}
		
		\pp
		
		\begin{ideaBox}
			Potencia elevada a potencia \(\rightarrow\) se multiplican exponentes.
		\end{ideaBox}
	\end{frame}
	
	% ----------------------------------------------------------
	\begin{frame}{\iconoDesarrollo\ \ Propiedad 4: potencia de un producto}
		\[
		(ab)^n=a^n b^n
		\]
		
		\pp
		
		\begin{exampleblock}{Ejemplo}
			\[
			(2\cdot5)^3=2^3\cdot5^3=8\cdot125=1000
			\]
		\end{exampleblock}
	\end{frame}
	
	% ----------------------------------------------------------
	\begin{frame}{\iconoDesarrollo\ \ Propiedad 5: potencia de un cociente}
		\[
		\left(\frac{a}{b}\right)^n=\frac{a^n}{b^n}, \qquad b\neq0
		\]
		
		\pp
		
		\begin{exampleblock}{Ejemplo}
			\[
			\left(\frac{2}{3}\right)^3=\frac{2^3}{3^3}=\frac{8}{27}
			\]
		\end{exampleblock}
	\end{frame}
	
	% ----------------------------------------------------------
	\begin{frame}{\iconoDesarrollo\ \ Propiedad 6: exponente cero}
		\[
		a^0=1,\qquad a\neq0
		\]
		
		\pp
		
		\begin{exampleblock}{Ejemplo}
			\[
			7^0=1
			\]
		\end{exampleblock}
	\end{frame}
	
	% ----------------------------------------------------------
	\begin{frame}{\iconoDesarrollo\ \ Propiedad 7: exponente negativo}
		\[
		a^{-n}=\frac{1}{a^n},\qquad a\neq0
		\]
		
		\pp
		
		\begin{exampleblock}{Ejemplo}
			\[
			2^{-3}=\frac{1}{2^3}=\frac{1}{8}
			\]
		\end{exampleblock}
	\end{frame}
	
	% ----------------------------------------------------------
	\begin{frame}{\iconoDesarrollo\ \ Propiedad 8: casos especiales}
		\[
		a^1=a
		\qquad
		1^n=1
		\qquad
		0^n=0 \text{ para } n>0
		\]
		
		\pp
		
		\begin{exampleblock}{Ejemplos}
			\[
			9^1=9
			\qquad
			1^{12}=1
			\qquad
			0^5=0
			\]
		\end{exampleblock}
	\end{frame}
	
	% ----------------------------------------------------------
	\begin{frame}{\iconoError\ \ Errores frecuentes}
		\begin{columns}[T,totalwidth=\textwidth]
			\begin{column}{0.48\textwidth}
				\begin{alertblock}{Error 1}
					\[
					a^m+a^n \neq a^{m+n}
					\]
					\[
					2^2+2^3=12
					\]
					\[
					2^5=32
					\]
				\end{alertblock}
			\end{column}
			\begin{column}{0.48\textwidth}
				\begin{alertblock}{Error 2}
					\[
					(a+b)^n \neq a^n+b^n
					\]
					\[
					(2+3)^2=25
					\]
					\[
					2^2+3^2=13
					\]
				\end{alertblock}
			\end{column}
		\end{columns}
	\end{frame}
	
	% ----------------------------------------------------------
	\begin{frame}{\iconoDesarrollo\ \ Patrones numéricos}
		\begin{columns}[T,totalwidth=\textwidth]
			\begin{column}{0.48\textwidth}
				\begin{block}{Potencias de 2}
					\[
					2^1=2,\quad 2^2=4,\quad 2^3=8
					\]
					\[
					2^4=16,\quad 2^5=32
					\]
				\end{block}
			\end{column}
			\begin{column}{0.48\textwidth}
				\begin{exampleblock}{Potencias de 10}
					\[
					10^1=10,\quad 10^2=100
					\]
					\[
					10^3=1000,\quad 10^4=10000
					\]
				\end{exampleblock}
			\end{column}
		\end{columns}
		
		\pp
		
		\begin{ideaBox}
			Si el exponente aumenta en 1, la potencia se multiplica por la base.
		\end{ideaBox}
	\end{frame}
	
	% ----------------------------------------------------------
	\SectionFrame{\iconoPractica}{Práctica guiada}{Resolución paso a paso}
	
	% ----------------------------------------------------------
	\begin{frame}{\iconoPractica\ \ Ejercicio guiado 1}
		\begin{block}{Problema}
			Simplificar:
			\[
			2^4\cdot2^3
			\]
		\end{block}
		
		\only<2>{
			\begin{exampleblock}{Paso 1}
				Misma base:
				\[
				2^4\cdot2^3=2^{4+3}
				\]
			\end{exampleblock}
		}
		
		\only<3>{
			\begin{exampleblock}{Paso 2}
				\[
				2^{4+3}=2^7=128
				\]
			\end{exampleblock}
		}
	\end{frame}
	
	% ----------------------------------------------------------
	\begin{frame}{\iconoPractica\ \ Ejercicio guiado 2}
		\begin{block}{Problema}
			Resolver:
			\[
			(3^2)^3
			\]
		\end{block}
		
		\only<2>{
			\begin{exampleblock}{Paso 1}
				Potencia de una potencia:
				\[
				(3^2)^3=3^{2\cdot3}
				\]
			\end{exampleblock}
		}
		
		\only<3>{
			\begin{exampleblock}{Paso 2}
				\[
				3^{2\cdot3}=3^6=729
				\]
			\end{exampleblock}
		}
	\end{frame}
	
	% ----------------------------------------------------------
	\SectionFrame{\iconoAgro}{Aplicación agronómica}{Matemática al servicio del contexto profesional}
	
	% ----------------------------------------------------------
	\begin{frame}{\iconoAgro\ \ Caso 1: semillas por superficie}
		\begin{block}{Situación}
			En una parcela experimental se usan:
			\[
			2^5
			\]
			semillas por metro cuadrado.
			
			La superficie cultivada es:
			\[
			2^3
			\]
			metros cuadrados.
		\end{block}
		
		\pp
		
		\begin{exampleblock}{Resolución}
			\[
			2^5\cdot2^3=2^{8}=256
			\]
			Se utilizan \textbf{256 semillas}.
		\end{exampleblock}
	\end{frame}
	
	% ----------------------------------------------------------
	\begin{frame}{\iconoAgro\ \ Caso 2: crecimiento bacteriano}
		\begin{block}{Situación}
			Una colonia bacteriana se duplica cada hora.
			
			Si parte con:
			\[
			2^3
			\]
			bacterias, luego de 4 horas habrá:
			\[
			2^3\cdot2^4
			\]
		\end{block}
		
		\pp
		
		\begin{exampleblock}{Resolución}
			\[
			2^3\cdot2^4=2^7=128
			\]
			Habrá \textbf{128 bacterias}.
		\end{exampleblock}
	\end{frame}
	
	% ----------------------------------------------------------
	\begin{frame}{\iconoAgro\ \ Caso 3: fertilización}
		\begin{block}{Situación}
			Se aplican:
			\[
			5\times10^2 \text{ gramos por hectárea}
			\]
			en una superficie de:
			\[
			10^3 \text{ hectáreas}
			\]
		\end{block}
		
		\pp
		
		\begin{exampleblock}{Resolución}
			\[
			(5\times10^2)(10^3)=5\times10^5=500000
			\]
			Se requieren \textbf{500000 gramos}.
		\end{exampleblock}
	\end{frame}
	
	% ----------------------------------------------------------
	\SectionFrame{\iconoPractica}{Práctica autónoma}{Consolidación individual}
	
	% ----------------------------------------------------------
	\begin{frame}{\iconoPractica\ \ Ejercicios propuestos}
		\begin{enumerate}
			\item \(2^5\cdot2^2\)
			\item \(\dfrac{3^7}{3^4}\)
			\item \((4^2)^3\)
			\item \((2\cdot3)^2\)
			\item \(\left(\dfrac{3}{5}\right)^2\)
			\item \(8^0\)
			\item \(2^{-4}\)
			\item En un cultivo hay \(3^4\) plantas por sector y \(3^2\) sectores. ¿Cuántas plantas hay en total?
		\end{enumerate}
	\end{frame}
	
	% ----------------------------------------------------------
	\begin{frame}{\iconoPractica\ \ Respuestas}
		\begin{multicols}{2}
			\begin{enumerate}
				\item \(2^5\cdot2^2=2^7=128\)
				\item \(\dfrac{3^7}{3^4}=3^3=27\)
				\item \((4^2)^3=4^6=4096\)
				\item \((2\cdot3)^2=36\)
				\item \(\left(\dfrac{3}{5}\right)^2=\dfrac{9}{25}\)
				\item \(8^0=1\)
				\item \(2^{-4}=\dfrac{1}{16}\)
				\item \(3^4\cdot3^2=3^6=729\)
			\end{enumerate}
		\end{multicols}
	\end{frame}
	
	% ----------------------------------------------------------
	\SectionFrame{\iconoCierre}{Cierre}{Síntesis y reflexión final}
	
	% ----------------------------------------------------------
	\begin{frame}{\iconoCierre\ \ Ideas clave}
		\begin{block}{Síntesis}
			\begin{itemize}
				\item<1-> Una potencia representa una multiplicación repetida.
				\item<2-> La base es el número que se repite.
				\item<3-> El exponente indica cuántas veces se repite la base.
				\item<4-> Las propiedades simplifican cálculos.
				\item<5-> Las potencias modelan crecimiento, concentración y escala.
			\end{itemize}
		\end{block}
	\end{frame}
	
	% ----------------------------------------------------------
	\begin{frame}{\iconoCierre\ \ Pregunta de salida}
		\begin{exampleblock}{Reflexión final}
			¿Por qué las potencias son una herramienta útil para describir fenómenos agrícolas como la siembra, la fertilización o el crecimiento bacteriano?
		\end{exampleblock}
		
		\vspace{0.5cm}
		
		\begin{center}
			{\Large\color{USTTeal}\faSeedling}
			
			\vspace{0.25cm}
			
			{\normalsize\textcolor{USTGray}{La matemática básica permite interpretar mejor la realidad técnica y profesional.}}
		\end{center}
	\end{frame}
	
\end{document}